
\documentclass[../../main.tex]{subfiles}
\begin{document}

% 年度の大きな表示
\begin{center}
  {\fontsize{48pt}{58pt}\selectfont \textbf{1987}\normalsize\textbf{年度}}
\end{center}

\vspace{10mm}

% 本文

\vspace{15mm}

% 出題テーマと難易度の表
\noindent\textbf{出題テーマと難易度}

\vspace{5mm}

\begin{center}
  \shadowbox{%
    \begin{tabular}{|c|p{8cm}|c|}
      \hline
      \textbf{問題番号} & \textbf{テーマ} & \textbf{難易度} \\
      \hline
      第1問 & 三次方程式の解の範囲 & \\
      \hline
      第2問 & 一次変換と数列の極限 & \\
      \hline
      第3問 & 円と三角形の外接円 & \\
      \hline
      第4問 & 媒介変数曲線と面積 & \\
      \hline
      第5問 & 確率（格子点上の経路） & \\
      \hline
    \end{tabular}%
  }
\end{center}

\vspace{15mm}

% 解答
\noindent\textbf{解答}

\vspace{5mm}

\begin{center}
  \shadowbox{%
    \renewcommand{\arraystretch}{1.6}
    \begin{tabular}{|c|c|p{8cm}|}
      \hline
      \textbf{問題番号} & \textbf{小問} & \textbf{答え} \\
      \hline
      第1問 & - & 証明問題 \\
      \hline
      第2問 & - &  \\
      \hline
      \multirow{2}{*}{第3問} & (1) & 証明問題 \\
      \cline{2-3}
                             & (2) & 図示（$|x|>1$ または $x^2+y^2<1$（$y\ne0$）かつ $x^2+y^2\le4$の範囲） \\
      \hline
      \multirow{2}{*}{第4問} & (1) & $a=\dfrac{1}{e}$ \\
      \cline{2-3}
                             & (2) & $e^2-2e$ \\
      \hline
      \multirow{2}{*}{第5問} & (1) & $P(k,l)={}_{k+l}C_k\left(\dfrac{1}{2}\right)^{k+l}$ \\
      \cline{2-3}
                             & (2) & $P(n,2)=\dfrac{n^2+5n+8}{8}\left(\dfrac{1}{2}\right)^n$ \\
      \hline
    \end{tabular}%
    \renewcommand{\arraystretch}{1.0}
  }
\end{center}

\vspace{15mm}

% 解答時間
\noindent\textbf{試験形態}

\vspace{5mm}

\begin{center}
  \shadowbox{%
    \begin{tabular}{p{8cm}}
      （要確認）
    \end{tabular}%
  }
\end{center}

\end{document}
